\documentclass[11pt,a4paper]{article}

% --- Paquetes Fundamentales ---
\usepackage[utf8]{inputenc}
\usepackage[T1]{fontenc}
\usepackage[spanish,es-tabla]{babel}
\usepackage[margin=2.5cm, headheight=40pt]{geometry}
\usepackage{graphicx}
\usepackage{fancyhdr}
\usepackage{xcolor}
\usepackage[colorlinks=true, linkcolor=blue, urlcolor=blue]{hyperref}

% --- Matemáticas y Electrónica ---
\usepackage{amsmath, amssymb}
\usepackage{siunitx}
\usepackage[american]{circuitikz}

% --- Elementos de Diseño ---
\usepackage{tcolorbox}
\usepackage{booktabs}
\usepackage{tabularx}
\usepackage{enumitem}
\usepackage{listings}

% --- Configuración de siunitx ---
\sisetup{
    output-decimal-marker = {,},
    per-mode = symbol,
    inter-unit-product = \ensuremath{{}\cdot{}}
}

% --- Configuración de Listings (Python) ---
\lstset{
    language=Python,
    basicstyle=\ttfamily\small,
    keywordstyle=\color{blue}\bfseries,
    commentstyle=\color{green!50!black}\itshape,
    stringstyle=\color{red!70!black},
    showstringspaces=false,
    frame=single,
    backgroundcolor=\color{gray!5},
    rulecolor=\color{gray!40},
    numbers=left,
    numberstyle=\tiny\color{gray}
}

% --- Estilo de Página ---
\pagestyle{fancy}
\fancyhf{}
\lhead{\textbf{GUÍA DE LABORATORIO N$^\circ$ 1} \\ IMT-121 Circuitos Electrónicos I}
\rhead{Carrera de Ingeniería Mecatrónica y Biomédica \\ Universidad Católica Boliviana ``San Pablo'' -- Tarija}
\cfoot{Página \thepage}
\renewcommand{\headrulewidth}{0.4pt}

\begin{document}

\begin{center}
    \Large\textbf{Formulación Matricial Canónica por Inspección Directa y Validación Tripartita de Redes Resistivas en CC}
\end{center}

\vspace{0.3cm}

\noindent
\textbf{Docente:} M.Sc. Ing. Bernardo Quiroga Turdera \\
\textbf{Gestión:} Semestre II/2026 \hspace{2cm} \textbf{Duración:} 120 minutos \\
\textbf{Modalidad:} Equipos de 3 estudiantes (Tríadas oficiales del PFI). \\
\textbf{Enfoque:} CDIO (Implementar--Operar) y Criterio ABET 6 (Experimentación y Análisis).

\vspace{0.4cm}
\hrule
\vspace{0.4cm}

\section{Objetivos de Aprendizaje}
\begin{itemize}
    \item Formular por inspección visual directa la matriz canónica $\mathbf{R}$ ($3\times 3$) y $\mathbf{G}$ para una red de CC.
    \item Implementar la solución computacional en Python (NumPy) y simulación \texttt{.op} en LTSpice XVII.
    \item Construir la red física midiendo experimentalmente voltajes nodales y corrientes de lazo.
    \item Validar la convergencia (teoría, simulación, hardware) mediante Error Relativo Porcentual ($\varepsilon_r \le 5\%$).
\end{itemize}

\section{Recursos y Equipamiento Requerido}
\textbf{a. Hardware y Componentes:}
\begin{itemize}[label=\raisebox{0.25ex}{\tiny$\bullet$}]
    \item Fuente de alimentación regulada de CC ($\pm\qty{12}{\volt}$ o $\qty{0}{\volt}$--\qty{30}{\volt}).
    \item Multímetro digital y Protoboard con jumpers.
    \item Resistencias (1/4 W): $R_1 = \qty{120}{\ohm}$, $R_2 = \qty{220}{\ohm}$, $R_3 = \qty{470}{\ohm}$, $R_4 = \qty{330}{\ohm}$, $R_L = \qty{100}{\ohm}$.
\end{itemize}

\textbf{b. Software Científico:} Python 3.12 (NumPy, Matplotlib) y LTSpice XVII.

\section{Desarrollo de la Práctica}

\subsection*{PRINCIPIO 1: Fundamentación Teórica (Previo de Laboratorio)}
\textit{Atención: Esta etapa debe ser desarrollada antes de ingresar al laboratorio.}

Dada la red de $3\times 3$ mallas (Red de Polarización de Sensor/Electrodo - Hito 1 del PFI):

\begin{center}
    \begin{circuitikz}[scale=1, transform shape]
        \draw 
        (0,0) to[V, v=$V_{s1}$] (0,3)
              to[R, l=$R_1$] (3,3)
              to[R, l=$R_2$, *-*] (3,0) -- (0,0);
        \draw 
        (3,3) to[R, l=$R_3$] (6,3)
              to[R, l=$R_4$, *-*] (6,0) -- (3,0);
        \draw 
        (6,3) to[R, l=$R_L$] (9,3)
              to[V, v_=$V_{s2}$] (9,0) -- (6,0);
    \end{circuitikz}
\end{center}

\begin{enumerate}
    \item Mida con el multímetro el valor real $R_{\text{medido}}$ de cada resistencia y registre tolerancias.
    \item Escriba la matriz canónica por inspección directa $\mathbf{R}\mathbf{I} = \mathbf{V}$:
    $$
    \begin{bmatrix} R_1 + R_2 & -R_2 & 0 \\ -R_2 & R_2 + R_3 + R_4 & -R_4 \\ 0 & -R_4 & R_4 + R_L \end{bmatrix} \begin{bmatrix} I_1 \\ I_2 \\ I_3 \end{bmatrix} = \begin{bmatrix} V_{s1} \\ 0 \\ -V_{s2} \end{bmatrix}
    $$
    \item Resuelva analíticamente $I_3$ y $V_{RL} = I_3 \cdot R_L$ usando la Regla de Cramer con valores medidos.
\end{enumerate}

\subsection*{PRINCIPIO 2: Implementación (Simulación + Código + Hardware)}
\textbf{A. Simulación en LTSpice XVII:} Dibuje el esquema, asigne nombres de red (\texttt{N\_nodo1}, \texttt{N\_salida}), configure \texttt{.op} y exporte resultados.\\
\textbf{B. Cómputo en Python:} Ejecute el siguiente bloque con sus valores medidos reales:
\begin{lstlisting}
import numpy as np
# 1. Valores reales medidos
R1, R2, R3, R4, RL = 121.2, 218.5, 465.0, 328.0, 99.1
Vs1, Vs2 = 15.00, -5.00

# 2. Formulacion Matricial
R = np.array([
    [R1 + R2, -R2, 0.0],
    [-R2, R2 + R3 + R4, -R4],
    [0.0, -R4, R4 + RL]
], dtype=float)
V = np.array([Vs1, 0.0, -Vs2], dtype=float)

# 3. Solucion
I_teorico = np.linalg.solve(R, V)
print("Corrientes (mA):", np.round(I_teorico * 1000, 4))
\end{lstlisting}
\textbf{C. Montaje Físico:} Ensamble en protoboard ($V_{s1}=\qty{15}{\volt}$, $V_{s2}=\qty{-5}{\volt}$). Mida voltajes nodales y corriente $I_L$.

\subsection*{PRINCIPIO 3: Validación Tripartita y Análisis de Error}
Fórmula de Error Relativo: $\varepsilon_r = \left\vert \frac{X_{\text{medido}} - X_{\text{teórico}}}{X_{\text{teórico}}} \right\vert \times 100\%$

\begin{table}[h!]
    \centering
    \begin{tabularx}{\textwidth}{@{}lXXXX@{}}
        \toprule
        \textbf{Variable} & \textbf{Teórico (Py)} & \textbf{Simulado (SPICE)} & \textbf{Medido (Real)} & \textbf{Error $\varepsilon_r$ (\%)} \\
        \midrule
        Voltaje Nodo 1 ($V_1$) & & & & \\
        Voltaje Carga ($V_{RL}$) & & & & \\
        Corriente Carga ($I_L$) & & & & \\
        Potencia Carga ($P_L$) & & & & \\
        \bottomrule
    \end{tabularx}
\end{table}

\textbf{Discusión Técnica (Informe):} ¿Justifica el Efecto de Carga del Multímetro ($R_{\text{int\_amp}} \approx \qty{1}{\ohm}$ y $R_{\text{vol}} \approx \qty{10}{\mega\ohm}$) las discrepancias en sus medidas?

\subsection*{PRINCIPIO 4: Evidencias y Entregables}
Subir a GitHub (\texttt{/Laboratorio\_1}): Script Jupyter (\texttt{.ipynb}), Archivos LTSpice (\texttt{.asc}, \texttt{.log}), Fotografía del montaje y \textbf{Informe Técnico IEEE (Max. 2 páginas)}.

\section{Rúbrica de Evaluación (100 Puntos)}
\begin{table}[h!]
    \centering
    \footnotesize
    \begin{tabularx}{\textwidth}{@{}p{2.5cm}XXXc@{}}
        \toprule
        \textbf{Criterio} & \textbf{Insuficiente (0-50\%)} & \textbf{Competente (51-85\%)} & \textbf{Sobresaliente (86-100\%)} & \textbf{Pond.} \\
        \midrule
        \textbf{Modelado y Cómputo} & Inhabilidad para armar la matriz $R$ o errores graves en NumPy. & Matriz $R$ correcta, con pequeños errores en la medición previa. & Formulación impecable. Script modular automatizado con valores reales. & 30\% \\
        \textbf{Simulación y Montaje} & Circuito mal armado o simulación inconclusa. & Circuito funcional, pero desordenado o sin calibración de fuentes. & Montaje profesional. Simulación \texttt{.op} con nodos claros y calibración precisa. & 30\% \\
        \textbf{Validación ($\varepsilon_r$)} & Omite la tabla o el cálculo de error porcentual. & Completa la tabla, pero no justifica causas físicas de las discrepancias. & Análisis riguroso. Discute efecto de carga de instrumentos. & 30\% \\
        \textbf{Reporte IEEE} & No sube archivos o entrega fuera de plazo. & Sube archivos, pero informe PDF no sigue plantilla IEEE. & Repositorio impecable. Informe IEEE sintético y profesional. & 10\% \\
        \bottomrule
    \end{tabularx}
\end{table}

\end{document}